\documentclass[10pt,a4paper]{amsart}
\usepackage[margin=19mm]{geometry}
\usepackage{amsmath,amssymb,mathtools}
\usepackage[scr=rsfs,cal=euler]{mathalfa}
\usepackage{xfrac}
\usepackage[unicode,hidelinks,bookmarks=false]{hyperref}
\newcommand{\ie}{i.e.\ }
\newcommand{\cf}{cf.\ }
\newcommand{\resp}{respectively\ }
\newcommand{\Subsection}{\subsection}
\providecommand{\nobreakdash}{\nobreak\hbox{-}}
\setlength{\emergencystretch}{2em}
\multlinegap0pt
\usepackage{bm}
\usepackage{bbm}
\usepackage{dsfont}
\usepackage{xspace}
\usepackage{cancel}
\usepackage{mathtools}
\usepackage{amsthm}
\makeatletter
\@ifundefined{thm}{\newtheorem{thm}{Thm}}{}
\@ifundefined{prop}{\newtheorem{prop}[thm]{Proposition}}{}
\makeatother
\newcommand{\Starphi}{S_\phi^*}
\title[Selberg's Central Limit Theorem weighted by Linear Statistic]{Selberg's Central Limit Theorem weighted by Linear Statistics of Zeta Zeros}
\author{Alessandro Fazzari, Maxim Gerspach, Paolo Minelli}
\date{Source: arXiv:2507.04150v1 (CC BY 4.0)}
\begin{document}
\maketitle

\noindent\emph{Source and licence.} Selberg's Central Limit Theorem weighted by Linear Statistics of Zeta Zeros. Version arXiv:2507.04150v1 (CC BY 4.0), distributed under the Creative Commons Attribution 4.0 International licence, as recorded in the arXiv licence metadata for this exact version and shown on the abstract page https://arxiv.org/abs/2507.04150v1. Full rights notice: licenses/arxiv-2507.04150-CC-BY-4.0.txt.\par\smallskip

\noindent\textit{Editorial scope.} Counted unit: the Prop whose source label is \texttt{prop:lkmoments} in the article (source0.tex lines 531--543, source label \texttt{prop:lkmoments}), delivered together with the article's own complete original proof of it (source0.tex lines 545--611, one \texttt{proof} environment of 67 lines). No mathematics is written, repaired or re-derived: statement and proof are the article's own text line for line. Required context retained uncounted: none. Every definition, display and label of those lines is retained unchanged. Omitted from this package and not redistributed: 1--530, 544--544, 612--727. Nothing inside a retained excerpt is omitted or altered: every retained body is delivered exactly as the article's lines are written, so no omission receipt of any kind accompanies this package. Citations: the retained excerpts carry no citation command at all, so no bibliography entry of the article is transcribed and no reference is redistributed. Reference adaptation: none was needed and none was made; every reference inside the delivered unit resolves inside it, so no unresolved reference or citation is produced. Printed slips kept literal: no printed word, symbol, hypothesis or number of the counted statement or of its proof was altered. Layout and class: the article's own class is \texttt{amsart} with packages including fontenc, inputenc, xcolor, datetime, natbib, comment, todonotes, bbm, amstext, amsthm, amssymb, xparse, mathrsfs, mathtools; the wrapper is the lane's pinned \texttt{amsart} editorial wrapper at 10pt on A4, which restates the theorem-like environments the retained text needs and adds the article's own macro definitions that the retained text needs (\texttt{\textbackslash Starphi}), with extra wrapper packages (bm, bbm, dsfont, xspace, cancel, mathtools). The delivered numbering is the unit's own. Assets: no retained span contains a figure environment, a graphics-inclusion command, a PDF-page inclusion, a table or a TikZ picture; no image file is copied into this package, the package declares no asset dependency and no image file is redistributed with it. Licence: version arXiv:2507.04150v1 of this article is distributed under the Creative Commons Attribution 4.0 International licence, as recorded in the arXiv licence metadata for this exact version and shown on the abstract page https://arxiv.org/abs/2507.04150v1. A full rights notice is delivered at licenses/arxiv-2507.04150-CC-BY-4.0.txt. Characters: every retained excerpt is pure ASCII, so the delivered engine, XeLaTeX with its default Latin Modern text font, reports no missing character.
\begin{prop}\label{prop:lkmoments}
    Let $T \ge 3$ be a (large) parameter and let $h,\ell,k$ be fixed nonnegative integers. Set $x = T^{\varepsilon}$ for some sufficiently small constant $\varepsilon = \varepsilon(h,\ell,k,\eta) > 0$ and define 
    \begin{align*}
           M(h,\ell,k) = M(h, \ell,k; T) := \frac{1}{T}\int_T^{2 T} P_x(t)^{h} \overline{P_x(t)}^{\ell} \Starphi(t)^k\, dt.
    \end{align*}
    Then, for some sufficiently small constant $\delta>0$, we have
    \begin{equation}\begin{split}\label{eq:MhkDiag}
        M(h,\ell, k) = \frac{(-1)^k}{(\log T)^k} 
        \sum_{J \subseteq \{1, \dots, k \}} \sum_{\substack{p_1, \dots, p_h\leq x\\
        q_1, \dots, q_\ell\leq x\\ n_1, \dots, n_k \ge 1 }} &\frac{\Lambda^*(n_1) \cdots \Lambda^*(n_k)}{\sqrt{p_1\cdots p_h q_1\dots q_\ell n_1 \cdots n_k}} \prod_{j=1}^k  \hat\phi\left(\frac{\log n_j}{\log T}\right) \times \\
        &\times\mathbbm{1} \bigg( \prod_{i\leq h} p_i \prod_{j \in J} n_j = \prod_{i\leq \ell} q_i \prod_{j\in J^c} n_j \bigg) + O (T^{-\delta}).
    \end{split}\end{equation}
\end{prop}

\begin{proof}
    By definition of $P_x$ and $\Starphi$, writing cosines as exponentials and multiplying out, we get
    \begin{equation}\begin{split}\label{eq:MhkDef}
        M(h,\ell, k) = \frac{(-1)^k}{(\log T)^k} \sum_{J \subseteq \{1, \dots, k \}} &\sum_{\substack{p_1, \dots, p_h  \leq x\\
        q_1, \dots, q_\ell\leq x\\
        n_1, \dots, n_k \ge 1}} \frac{\Lambda^*(n_1) \cdots \Lambda^*(n_k)}{\sqrt{p_1\cdots p_h q_1\dots q_\ell n_1 \cdots n_k }} \times \\ \times &\prod_{j=1}^k \hat\phi\left(\frac{\log n_j}{\log T}\right) \frac{1}{T} \int_T^{2 T} \left( \frac{\prod_{i\leq \ell} q_i \prod_{j \in J^c} n_j}{\prod_{i\leq h} p_i \prod_{j \in J} n_j} \right)^{i t} \, dt.
    \end{split}\end{equation}
    To deduce the claim, we first recall the standard bound
    \begin{equation}\label{eq:xitBound}
        \int_T^{2T} x^{it} \, dt \ll \frac{1}{| \log x |}
    \end{equation} 
    for $x \ne 1$, which follows immediately from integrating.
    Let us fix $J\subseteq\{1, \dots, k\}$  and denote $$P(J):=\frac{\prod_{i\leq \ell} q_i \prod_{j \in J^c} n_j}{\prod_{i\leq h} p_i \prod_{j \in J} n_j}.$$ 
    We split the range of summation in \eqref{eq:MhkDef} according to whether
    \begin{equation}\label{eq:QuotNear1}
        \tfrac{1}{2} \le P(J)  \le 2
    \end{equation}
    holds. 
    
    If the inequality \eqref{eq:QuotNear1} does not hold, then by \eqref{eq:xitBound} the integral appearing in \eqref{eq:MhkDef} is bounded. The contribution from such tuples to \eqref{eq:MhkDef} is thus
    \begin{align*}
        \ll \frac{1}{T} \sum_{J \subseteq \{1, \dots, k \}} \sum_{\substack{p_1, \dots, p_h\leq x \\
        q_1, \dots, q_\ell\leq x\\
        n_1, \dots, n_k \le T^\eta
        }} \frac{1}{\sqrt{p_1\cdots p_h  q_1\dots q_\ell n_1 \cdots n_k}} 
        \ll \frac{x^{(h+\ell)/2}T^{\eta k/2}}{T}
        \ll T^{-\delta},
    \end{align*}
    for some fixed small $\delta>0$.
    We can hence restrict the sum in \eqref{eq:MhkDef} to tuples satisfying \eqref{eq:QuotNear1}, up to an error of $\ll T^{-\delta}$. Clearly, the case $P(J)=1$
    corresponds to the main term in \eqref{eq:MhkDiag}. 
    
    For the remaining range $P(J)\in[\frac12,2] \setminus \{ 1 \}$, employing \eqref{eq:xitBound}, it thus suffices to show that, for any $J \subseteq \{1,\dots,k\}$,
    \begin{equation}\label{eq:QuotNear1STS} 
        \sum_{\substack{p_1, \dots, p_h  \leq x\\ q_1, \dots, q_\ell\leq x \\ n_1, \dots, n_k \le T^\eta}} 
        \frac{1}{\sqrt{p_1\dots p_h q_1\dots q_\ell n_1\dots n_k}} \frac{\mathbbm{1} \left( P(J) \in \left[ \tfrac{1}{2}, 2 \right] \setminus \{1 \}  \right)}{\left| \log P(J) \right|}  
        \ll T^{1-\delta}.
    \end{equation}
    Note that we may assume without loss of generality that $|J| \ge \frac{k}{2}$; otherwise, one can simply exchange $J$ with $J^c$ and swap the roles of $q_i$ and $p_i$ in the following argument. 
    Now, writing
    \begin{align}\label{eq:shiftbyh}
        \prod_{i\leq h} p_i \prod_{j\in J} n_j=\prod_{i\leq \ell} q_i \prod_{j\in J^c} n_j +h
    \end{align}
    for some nonzero integer $h$ with $-T\leq h\leq T$, we have
    \begin{align*}
        \frac{1}{\vert \log P(J)\vert}\ll \frac{\prod_{i\leq h} p_i \prod_{j\in J} n_j}{\vert h\vert}.
    \end{align*}
    As a consequence, we deduce
    \begin{align*}
        \frac{1}{\sqrt{p_1\dots p_h q_1\dots q_\ell n_1\dots n_k}} \frac{1}{\left| \log P(J) \right|} 
        \ll  \frac{\prod_{j\leq h} p_i \prod_{j\in J} n_j}{|h|\sqrt{p_1\dots p_h q_1\dots q_\ell n_1\dots n_k}} 
        = \frac{1}{|h|\sqrt{P(J)}} 
        \ll \frac{1}{|h|}  ,
    \end{align*}
    by \eqref{eq:QuotNear1}. Thus, we have that the left-hand side of \eqref{eq:QuotNear1STS} is bounded by
    \begin{align*}
        &\ll \sum_{0\neq |h| \leq T} \frac{1}{|h|} \sum_{\substack{p_1, \dots, p_h  \leq x\\ q_1, \dots, q_\ell\leq x \\ n_1, \dots, n_k  \leq T^\eta}} \mathbbm{1} \left( \prod_{i\leq h} p_i \prod_{j\in J} n_j = \prod_{i\leq \ell} q_i \prod_{j\in J^c} n_j+h \right)\\
        &\ll \sum_{\substack{q_1,\dots,q_\ell \leq x \\ n_j\leq T^\eta,\, j\in J^c}}  \sum_{0\neq |h| \leq T} \frac{1}{|h|} \sum_{\substack{p_1, \dots, p_h  \leq x \\ n_j \leq T^\eta, \, j\in J}} \mathbbm{1} \left( \prod_{i\leq h} p_i \prod_{j\in J} n_j = \prod_{i\leq \ell} q_i \prod_{j\in J^c} n_j+h \right).
    \end{align*}
    Now, for any given $(q_i)_{i\leq \ell}$, $(n_j)_{j\in J^c}$, and $h$, the number of possible $(p_i)_{i\leq h},(n_j)_{j\in J}$ such that the indicator function above equals 1 is certainly $\ll T^{\varepsilon}$, say. Thus, the last line in the above display is bounded by
    \begin{align*}
        &\ll T^{\varepsilon} \sum_{\substack{q_1,\dots,q_\ell \leq x \\ n_j\leq T^\eta,\, j\in J^c}}  \sum_{0\neq |h| \leq T} \frac{1}{|h|}
        \ll T^{2\varepsilon}  x^{\ell} T^{\eta |J^c|} 
        \ll T^{(\ell + 2)\varepsilon} T^{\eta |J^c|}.
    \end{align*}
    Since $|J^c|\leq \frac k2$ by assumption, taking $\varepsilon$ small enough, we find that the above is $\ll T^{1-\delta}$. This concludes the proof of \eqref{eq:QuotNear1STS}, and therefore that of \eqref{eq:MhkDiag}.
\end{proof}

\end{document}
